\chapter{Efficient Inference}
\label{ch:inference}

BabyGPT's \code{generate()} method recomputes the full attention matrix at every step.  For a 1\,024-context model with 12 layers, each generated token touches roughly 100 million multiply-adds worth of redundant key/value projections.  In this chapter we implement a \textbf{KV cache} that stores previously computed keys and values, reducing per-token cost from $O(T)$ to $O(1)$ in the attention computation.  We also cover nucleus (top-$p$) sampling.


\section{The Autoregressive Bottleneck}

Generating $N$ tokens requires $N$ forward passes.  In each pass, self-attention computes keys and values for \emph{all} positions in the sequence --- including positions that were already processed.

Without a cache, generating token $t$ requires:
\begin{itemize}
  \item Computing Q, K, V for all $t$ positions: $O(t \cdot d^2)$
  \item Attention matrix: $O(t^2 \cdot d)$
\end{itemize}

The K and V for positions $0$ through $t{-}1$ are identical to what was computed when generating token $t{-}1$.  Recomputing them is pure waste.


\section{The KV Cache}

The fix: cache K and V for all previous positions.  When generating token $t$:
\begin{enumerate}
  \item Compute Q, K, V for position $t$ only ($O(d^2)$).
  \item Append K and V to the cache.
  \item Attend using Q for position $t$ against \emph{cached} K, V for all $t$ positions ($O(t \cdot d)$).
\end{enumerate}

Memory cost per layer:
\[
  2 \times T_{\max} \times n_{\text{kv\_head}} \times d_{\text{head}} \times \text{dtype\_size}
\]

For our model ($T{=}1024$, 12 KV heads, $d{=}64$, bf16):
\[
  2 \times 1024 \times 12 \times 64 \times 2 = 3.1\,\text{MB per layer} \times 12 = 37.7\,\text{MB total}
\]

With GQA (4 KV heads instead of 12), this drops to 12.6\,MB --- small enough to be irrelevant.

\subsection{Implementation}

\filetag{microgpt/generate.py}
\begin{lstlisting}
class KVCache:
    """Pre-allocated key/value cache for autoregressive generation.

    For each layer, stores keys and values of all previously processed
    tokens.  Shape per layer: (batch, max_seq_len, n_kv_head, head_dim).
    """

    def __init__(
        self,
        n_layers: int,
        max_seq_len: int,
        n_kv_head: int,
        head_dim: int,
        batch_size: int = 1,
        dtype: torch.dtype = torch.bfloat16,
        device: str = "cuda",
    ) -> None:
        self.n_layers = n_layers
        self.max_seq_len = max_seq_len
        self.pos = 0  # next write position

        # Pre-allocate: (n_layers, 2, batch, max_seq_len, n_kv_head, head_dim)
        # The '2' dimension is for K and V
        self.cache = torch.zeros(
            n_layers, 2, batch_size, max_seq_len, n_kv_head, head_dim,
            dtype=dtype, device=device,
        )

    def update(
        self,
        layer_idx: int,
        k: torch.Tensor,
        v: torch.Tensor,
    ) -> tuple[torch.Tensor, torch.Tensor]:
        """Write new K, V and return the full cached K, V.

        k, v: (B, T_new, n_kv_head, head_dim) --- new tokens
        Returns: (B, T_total, n_kv_head, head_dim) --- all cached tokens
        """
        T_new = k.size(1)
        self.cache[layer_idx, 0, :, self.pos : self.pos + T_new] = k
        self.cache[layer_idx, 1, :, self.pos : self.pos + T_new] = v
        # Return everything up to and including the new tokens
        end = self.pos + T_new
        return (
            self.cache[layer_idx, 0, :, :end],
            self.cache[layer_idx, 1, :, :end],
        )

    def advance(self, n: int) -> None:
        """Advance the write position by n tokens."""
        self.pos += n
\end{lstlisting}


\section{Prefill and Decode Phases}

Generation has two phases:

\textbf{Prefill}: process the entire prompt at once.  This fills the KV cache for all prompt positions.  Attention is causal over the prompt length.

\textbf{Decode}: generate one token at a time.  Each step computes Q for the new position only and attends against the full cached K, V.  Attention is not causal during decode (the single query can attend to all cached positions).

\tip{The prefill phase can be batched and parallelized just like training --- it is compute-bound.  The decode phase is memory-bound because each step only processes one token but reads the full KV cache.  This is why KV cache size matters for inference throughput.}


\section{Top-$p$ (Nucleus) Sampling}

In addition to temperature and top-$k$ (covered in the BabyGPT tutorial), we implement \textbf{nucleus sampling} (Holtzman et al., 2020):

\begin{enumerate}
  \item Sort tokens by probability (descending).
  \item Compute the cumulative probability.
  \item Keep the smallest set of tokens whose cumulative probability $\geq p$.
  \item Re-normalize and sample.
\end{enumerate}

Top-$p$ is adaptive: when the model is confident, it keeps few tokens; when uncertain, it keeps many.  Typical values: $p = 0.9$ or $p = 0.95$.

\begin{lstlisting}[language=Python]
if top_p is not None:
    sorted_logits, sorted_indices = torch.sort(logits, descending=True)
    cumprobs = torch.cumsum(F.softmax(sorted_logits, dim=-1), dim=-1)
    remove_mask = cumprobs - F.softmax(sorted_logits, dim=-1) >= top_p
    sorted_logits[remove_mask] = float("-inf")
    logits = sorted_logits.scatter(1, sorted_indices, sorted_logits)
\end{lstlisting}


\section{Summary}

Efficient inference requires:
\begin{enumerate}
  \item \textbf{KV cache}: store K, V for previously processed tokens; compute only Q for new tokens.
  \item \textbf{Prefill + decode}: batch-process the prompt, then generate one token at a time.
  \item \textbf{Sampling strategies}: temperature, top-$k$, and top-$p$ for controlling output diversity.
\end{enumerate}

The full generation script is in \code{microgpt/generate.py}, including a CLI for generating text from a trained checkpoint.
